\documentclass[10pt]{article}

\usepackage[utf8]{inputenc}

\usepackage[T1]{fontenc}

\usepackage{amsmath}

\usepackage{amsfonts}

\usepackage{amssymb}

\usepackage[version=4]{mhchem}

\usepackage{stmaryrd}

\usepackage{bbold}



\begin{document}

компактного многообразия $M^{n}$ заменить на разложение с ровно одной ручкой индекса 0 . Если $M^{n}$ односвязно и $n>5$, то, складывая и сокращая ручки, можно любое разложение на ручки свести к разложению с минимальным числом ручек, совместимым с гомологич. структурой $M^{n}$.



Пусть $(H, h)$ - топологич. ручка индекса $k$ в кусочно линейном многообразии $\boldsymbol{M}^{n}$, причем характеристич. отображение $h$ кусочно линейно в окрестности $\partial B^{k} \times$ $\times B^{n-k}$. Существует ли неподвижная на окрестности $\partial B^{k} \times B^{n-k}$ изотопия отображения $h$, выпрямляющая ручку $H$, т. е. переводящая ее в кусочно линейную ручку? Если бы ответ на этот вопрос был всегда положительным, то на каждом топологич. многообразии можно было бы ввести кусочно линейную структуру, согласовывая структуры на координатных окрестностях с помощью выпрямления кусочно линейных ручек одной окрестности внутри другой. На самом деле ответ зависит от индекса $k$ и размерности $n$ ручки $H$. Если $n \leqslant 3$ или $n \geqslant 5$ и $k \neq 3$, то любая ручка выпрямляема. Известно, что при $n \geqslant 5$ существуют невыпрямляемые ручки индекса 3 , причем препятствие к выпрямлению лежит в группе $\mathbb{Z}_{2}$. В размерности 4 ручки индексов 0 и 1 выпрямляемы, а индексов 2 и 3 , вообще говоря, нет. Препятствием к сглаживанию кусочно линейных ручек служат т. н. группы Милнора $\Gamma_{k}$.



Лит.: [1] С м е й л С., “Математика", 1962, т. 6, №3, с. $138-55 ;[2]$ е г о же, там же, 1964, т. 8, № 4, с. 95-108; [4] K i r b y R. C., S i e b e n m a n n L. C., "Ann. Math. Stud.», 1977, № 88;'[5] Р у р к К., С а н д е р с о н Б., Введение в кусочно линейную топологию, пер. с англ., М., 1974; [6] Р о хл и н В. А., $\mathbf{9}$ у к с Д. Б., Начальный курс топологии. Геометрические главы, М., $1977 . \quad$ С. В. Матвеев. РУЧНОЕ ВЛОЖЕНИЕ - вложение топологич. полиәдра $P$ в пространство $\mathbb{R} n$ такое, что существует гомеоморфизм $\mathbb{R}^{n}$ на себя, при к-ром $\boldsymbol{P}$ переходит в прямолинейный полиэдр. Тогда $P$ наз. р уч ным. В противном случае $P$ наз. д и к и м, а вложение диким вложенией.



PУIII TEOPEMA: пусть $f(z)$ и $g(z)$ - регулярные аналитич. функции комплексного переменного $z$ в области $D$, простая замкнутая кусочно гладкая кривая $\Gamma$ вместе с ограничиваемой ею областью $G$ принадлежит $D$ и всюду на $\Gamma$ выполняется неравенство $|f(z)|>|g(z)|$; тогда в области $G$ сумма $f(z)+g(z)$ имеет столько же нулей, сколько и $f(z)$.



Эта теорема была получена Э. Руше [1]. Она является следствием аргумента принципа, и из нее в свою очередь получается основная теорема алгебры многочленов.



Справедливо также о б о б ш е н и е Р. т. для многомерных голоморфных отображений, напр. в следующем виде. Пусть $f(z)=\left(f_{1}(z), \ldots, f_{n}(z)\right)$ и $g(z)=\left(g_{1}(z), \ldots\right.$, $\left.g_{n}(z)\right)$ - голоморфные отображения области $D$ комплексного пространства $\mathbb{C}^{n}$ в $\mathbb{C}^{n}, n \geqslant 1$, с изолированными нулями, пусть гомеоморфная сфере гладкая поверхность $\Gamma$ вместе с ограничиваемой ею областью $G$ принадлежит $D$ и всюду на $\Gamma$ выполняется неравенство



\[

|f(z)|=\sqrt{\left|f_{1}(z)\right|^{2}+\ldots+\left|f_{n}(z)\right|^{2}}>|g(z)| .

\]



Тогда отображение $f(z)+g(z)$ имеет в $G$ столько же нулей, сколько и $f(z)$.



Jum.: [1] R o u c hé E., «J. École polyt.», 1858, t. 21; [2] М а p к у ш е в и ч А. И., Теория аналитических функций, 2

сизд., т. 1, М., 1967; [3] Шіа а б а т Б. В., Введение в комплек-

Б. Д. Соломенцев.



РЭЛЕЯ РАСІРЕДЕЛЕНИЕ - непрерывное распределение вероятностей с плотностью



\[

p(x)= \begin{cases}\frac{x}{\sigma^{2}} e^{-x^{2} / 2 \sigma^{2}}, & x>0, \\ 0 & , x \leqslant 0,\end{cases}

\]



зависящей от масштабного параметра $\boldsymbol{\sigma}>0$. Р. p. имеет положительную асимметрию, его единственная мода находится в точке $x=\sigma$. Все моменты Р. р. конечны, математич. ожидание и дисперсия равны соответственно $\sqrt{\frac{\pi}{2}} \sigma$ и $2 \sigma^{2}\left(1-\frac{\pi}{4}\right)$. Функция распределения Р. р. имеет вид



\[

F(x)= \begin{cases}1-e^{-x^{2} / 2 \sigma^{2}}, & x>0, \\ 0 & , x \leqslant 0 .\end{cases}

\]



P. р. является частным случаем распределения с плотностью



\[

\frac{2}{2^{n / 2} \sigma^{n} \Gamma\left(\frac{n}{2}\right)} x^{n-1} e^{-x^{2} / 2 \sigma^{2}},

\]



при $n=2$ и, следовательно, при $\sigma=1$ Р. р. совпадает с распределением арифметического квадратного корня из случайной величины, имеющей «xu-квадрат» распредeление с двумя степенями свободы. Иначе, Р. р. может быть интерпретировано как распределение длины вектора в прямоугольной системе координат на плоскости, координаты к-рого независимы и имеют норлальное распределение с параметрами 0 и $\sigma^{2}$. Аналогом Р. p. в 3-мерном пространстве служит Максвелла распределениe.



P. р. находит основное применение в теории стрельбы и статистич. теории связи. Р. р. впервые рассмотрено Рәлеем (Rayleigh, 1880), как распределение результирующей амплитуды при сложении гармонич. колебаний.



Лит.: [1] С т р е т т Д ж. В. (л о р Д Р ә л е й), Волновая теория света, пер. с англ., М.- Л., 1940. $\quad$ А. В. Прохоров. РӘЛЕЯ УРАВНЕНИЕ - нелинейное обыкновенное дифференциальное уравнение 2-го порядка



\[

\ddot{x}+F(\dot{x})+x=0, \quad \dot{x}=d x / d t,

\]



где функция $F(u)$ удовлетворяет предположению:



$u F(u)<0$ при малых $|u|$,



$u F(u)>0$ при больших $|u|$.



P. у. описывает типичную нелинейную систему с одной степенью свободы, в к-рой возможны автоколебания. Названо по имени Рәлея (Rayleigh), изучавшего уравнение такого типа в связи с задачами акустики [1].



Если уравнение (*) продиффференцировать, а затем положить $y=\dot{x}$, то получится Льенара уравнение



\[

\ddot{y}+f(y) \dot{y}+y=0 ; f(u)=F^{\prime}(u)

\]



Частным случаем Р. у. при



\[

F(u)=-\lambda\left(u-\frac{u^{3}}{3}\right), \lambda=\text { const, }

\]



является Ван дер Поля уравнение. Иногда Р. у. наз. частный случай уравнения $(*)$ :



\[

\ddot{x}-\left(a-b \dot{x}^{2}\right) \dot{x}+x=0, a>0, b>0 .

\]



Имеется большое число работ, в к-рых выясняются условия существования и единственности устойчивого предельного цикла у Р. у., то есть условия возникновения автоколебаний. Вопрос о периодич. решениях изучался и для различных обобщений P. у., напр. для



\[

\ddot{x}+F(x, \dot{x}) \dot{x}+g(x)=e(t)

\]



где $e(t)$ - периодич. Функция.



Системой ти п а Рәле я часто наз. уравнение



\[

\begin{gathered}

\ddot{x}+F(\dot{x})+G(x)=H(t, x, \dot{x}), \\

x \in \mathbb{R} n, F: \mathbb{R} n \longrightarrow \mathbb{R}^{n}, G \cdot \mathbb{R}^{n} \longrightarrow \mathbb{R}^{n},

\end{gathered}

\]





\end{document}